%
% 8-grenzen.tex -- Grenzen des Modells und aktive Dämpfung
%
% (c) 2026 Stephanie Märklin, OST Ostschweizer Fachhochschule
%
% !TEX root = ../../buch.tex
% !TEX encoding = UTF-8
%
% Rolle im Paper: Gültigkeitsbereich
% Roter Faden: Nur kleine Störungen, ein Flugzustand pro A,
% ausgeblendete Kopplung zur Spiralmode, Yaw Damper.
% Überleitung aus dem vorangehenden Abschnitt: Gilt das uneingeschränkt?
%
\section{Grenzen des Modells und aktive Dämpfung
\label{aircraft:section:grenzen}}
\kopfrechts{Grenzen des Modells}
Das Bewegungsgesetz gilt nach Abschnitt~\ref{aircraft:section:system} nur für
kleine Störungen um die Ruhelage.
Bei grossen Auslenkungen hängen die Kräfte und Momente nicht mehr linear von den
Zustandsgrössen ab, und die Derivative sind keine Konstanten mehr.
Bei hohen Anstellwinkeln kann die Rolldämpfung $L_p$ ihr Vorzeichen wechseln und
positiv werden \cite{aircraft:stengel}.
Die Rollbewegung, die nach Abschnitt~\ref{aircraft:section:1-grundbegriffe} bei
gewöhnlichen Anstellwinkeln rasch abklingt, wächst dann an.
Das lineare Modell sagt dieses Verhalten nicht voraus, weil es mit den
Derivativen der Ruhelage rechnet.

Die Werte in der Systemmatrix gelten für ein Flugzeug in einem Flugzustand.
Eine andere Geschwindigkeit, eine andere Höhe oder eine andere Klappenstellung
ergibt andere Derivative und damit andere Eigenwerte.
Ein Flugzeug muss das Stabilitätskriterium in jedem Flugzustand erfüllen, den
es erreicht.
Die Rechnung aus Abschnitt~\ref{aircraft:section:4-eigenwerte} ist deshalb für
jeden Flugzustand zu wiederholen.
Dabei kann sich auch die Art der Eigenwerte ändern.
Der Fall zweier komplexer Eigenwertpaare tritt ein, wenn Rollbewegung und
Spiralbewegung zu einer gemeinsamen Schwingung verschmelzen, wie
Stengel \cite{aircraft:stengel} am Versuchsflugzeug M2-F2 zeigt.

Wir haben uns in Abschnitt~\ref{aircraft:section:system} auf die
Seitenbewegung beschränkt und die Längsbewegung ausgeblendet.
Für kleine Störungen ist das gerechtfertigt, weil die beiden Bewegungen dann
voneinander unabhängig sind.
Bei schnellen Rollmanövern und bei hohen Anstellwinkeln koppeln Trägheitskräfte
und aerodynamische Kräfte die Seitenbewegung mit der Längsbewegung
\cite{aircraft:nelson}.
Diese Fälle liegen ausserhalb des Modells.

Liefert die Geometrie nicht genügend Dämpfung, so wird die Dämpfung künstlich
erzeugt.
Ein \emph{Gierdämpfer} (yaw damper) ist ein automatisches System, das die
Gierrate $r$ misst und das Seitenruder proportional dazu ausschlägt, so dass
das Ruder der Drehung entgegenwirkt \cite{aircraft:nelson}.
Das zusätzliche Giermoment ist proportional zu $r$ und wirkt in der
Systemmatrix wie eine grössere Gierdämpfung $N_r$.
Nach Abschnitt~\ref{aircraft:section:fluegel} verschiebt ein grösserer Betrag
von $N_r$ das Eigenwertpaar der Dutch Roll nach links.
In der Sprache von Abschnitt~\ref{aircraft:subsection:realteil} wird der
Realteil $\sigma$ kleiner, und die Störung klingt schneller ab.
Diese Rückkopplung einer gemessenen Grösse auf ein Ruder heisst
\emph{Stabilitätsaugmentation} (stability augmentation) und verändert die
Einträge der Systemmatrix durch Regelung \cite{aircraft:nelson}.
Die meisten Verkehrsflugzeuge sind mit einem Gierdämpfer ausgerüstet
\cite{aircraft:nelson}.
Ein Verkehrsflugzeug muss auch bei Ausfall des Gierdämpfers fliegbar bleiben,
und der Konstrukteur berechnet deshalb die Eigenwerte der ungeregelten und der
geregelten Systemmatrix.

Bei Kampfflugzeugen wird die Stabilität zugunsten der Wendigkeit meist
deutlich reduziert.
\cite{aircraft:sadraey}.
Für ein solches Flugzeug ist die Stabilitätsaugmentation keine Verbesserung,
sondern die Bedingung dafür, dass es fliegbar ist.
Ein \emph{automatisches Flugregelsystem} (automatic flight control system)
steuert die Ruder durchgehend.
Der Gierdämpfer verändert einen Eintrag der Systemmatrix, das
automatische Flugregelsystem verändert mehrere Einträge zugleich.
Über die Stabilität eines solchen Flugzeugs entscheiden die Eigenwerte der
geregelten Systemmatrix.